% ==============================================================================
% Chapitre 5 : Gouvernance et Gestion de Projet
% Organisation équipe, méthodologie, risques, plan de déploiement, jalons
% ==============================================================================

\chapter{Gouvernance et Gestion de Projet}

\section{Organisation de l'\'equipe}

L'\'equipe du projet LockBits est compos\'ee de quatre \'etudiants de
l'\textbf{ESGI}, chacun occupant un r\^ole sp\'ecifique qui refl\`ete \`a
la fois les comp\'etences individuelles et les besoins du projet. La
r\'epartition des responsabilit\'es suit une logique de compl\'ementarit\'e~:
les domaines techniques sont attribu\'es selon l'expertise de chacun,
tandis que la coordination g\'en\'erale et l'infrastructure DevOps sont
confi\'ees \`a un membre unique pour garantir la coh\'erence de
l'ensemble.

Cette organisation vise \`a maximiser l'autonomie de chaque d\'eveloppeur
sur son p\'erim\`etre tout en maintenant une vision d'ensemble partag\'ee
par l'\'equipe. Les choix architecturaux et les points de convergence
entre les composants sont discut\'es coll\'egialement lors des r\'eunions
hebdomadaires.

\subsection{R\^oles et responsabilit\'es}

Le tableau~\ref{tab:rôles} pr\'esente la r\'epartition des r\^oles au
sein de l'\'equipe.

\begin{table}[h]
  \centering
  \caption{R\'epartition des r\^oles et responsabilit\'es}
  \label{tab:rôles}
  \begin{tabularx}{\textwidth}{@{}l l X@{}}
    \toprule
    \textbf{R\^ole} & \textbf{Personne} & \textbf{Responsabilit\'es} \\
    \midrule
    Chef de Projet / DevOps / H\'ebergeur
      & Cl\'ery A.-Ferradou
      & Coordination \'equipe, infrastructure Proxmox,
        h\'ebergement LXC (Dokploy, cloudflared),
        CI/CD, Docker, monitoring, GitOps,
        d\'eploiement, Cloudflare Zero Trust,
        gestion des risques \\
    D\'eveloppeur EDR
      & Swann Kechar
      & Serveur FastAPI, agent Python, FIM, int\'egration
        VirusTotal/GLPI, compilation PyInstaller \\
    D\'eveloppeur Site
      & Soltane Afellah
      & Site PHP/Apache, MySQL, GLPI SSO, dashboard client \\
    D\'eveloppeur Chatbot
      & Quentin Labourdette
      & Int\'egration API Claude, interface chat,
        rate limiting \\
    \bottomrule
  \end{tabularx}
\end{table}

\subsection{R\'epartition des composants}

Chaque membre de l'\'equipe est responsable d'un ensemble de composants
et de fichiers sp\'ecifiques, d\'efinis dans les d\'ep\^ots Git
correspondants.

\textbf{Cl\'ery A.-Ferradou}~: Chef de Projet, DevOps et H\'ebergeur, il assure la
coh\'erence de l'infrastructure. Il g\`ere un serveur Proxmox personnel
qui h\'eberge l'ensemble des conteneurs LXC du projet. Son p\'erim\`etre couvre
l'orchestration Docker (fichiers \texttt{docker-compose.yml},
\texttt{Dockerfiles}), la pipeline CI/CD (fichier
\texttt{build-and-push.yml} GitHub Actions), le monitoring
(Prometheus, Grafana, Loki) et les scripts de d\'eploiement
(\texttt{install.sh}, \texttt{backup.sh}). Il a mis en place et administre
le LXC \textbf{Dokploy} pour le d\'eploiement automatis\'e ainsi que
le LXC \textbf{cloudflared} pour le tunnel Cloudflare Zero Trust.
L'ensemble des services est accessible publiquement via le domaine
\texttt{lockbits.pro} gr\^ace \`a cette infrastructure. Il coordonne \'egalement
les phases de d\'eveloppement et la gestion des risques.

\textbf{Swann Kechar}~: D\'eveloppeur EDR, il travaille sur le cœur
du produit. Il est responsable du serveur FastAPI (dossier
\texttt{edr/server/}), de l'agent Python (dossier \texttt{edr/agent/}),
de la surveillance des fichiers (FIM), de l'int\'egration avec
VirusTotal et GLPI, et de l'empaquetage via PyInstaller.

\textbf{Soltane Afellah}~: D\'eveloppeur web, il con\c{c}oit et
d\'eveloppe l'ensemble du site client (dossier
\texttt{site\_lockbits/}) en PHP, HTML, CSS et JavaScript, incluant le
sch\'ema de base de donn\'ees MySQL et l'int\'egration du SSO GLPI.

\textbf{Quentin Labourdette}~: D\'eveloppeur chatbot, il impl\'emente
l'assistant virtuel int\'egr\'e au site. Son travail porte sur le
fichier \texttt{chat.php}, l'interface de chat, l'API Claude
d'Anthropic et le m\'ecanisme de rate limiting.

\section{M\'ethodologie de travail}

\subsection{Points hebdomadaires}

L'\'equipe a adopt\'e une organisation l\'eg\`ere adapt\'ee au contexte
acad\'emique et \`a la taille r\'eduite de l'\'equipe. Plut\^ot qu'un
cadre Agile formel, le suivi repose sur des \textbf{r\'eunions
hebdomadaires} de coordination.

Chaque semaine, l'\'equipe se retrouve pour un point o\`u chaque membre
pr\'esente~:

\begin{itemize}
  \item Les t\^aches r\'ealis\'ees depuis la derni\`ere r\'eunion~;
  \item Les \'el\'ements bloquants \'eventuels~;
  \item Les travaux en cours~;
  \item Les prochaines \'etapes pr\'evues.
\end{itemize}

Ce format l\'eg-er permet de maintenir une dynamique d'\'equipe sans
alourdir l'emploi du temps, tout en assurant une visibilit\'e partag\'ee
sur l'avancement de chacun. Les t\^aches sont suivies dans un tableau
Kanban GitHub Projects, offrant une vue d'ensemble accessible \`a tout
moment.

\subsection{Rituels}

Les rituels de l'\'equipe sont organis\'es autour des r\'eunions
hebdomadaires qui structurent la vie du projet.

\begin{description}
  \item[R\'eunion hebdomadaire]~: point de coordination d'environ 15 \`a
    30 minutes. Chaque membre pr\'esente tour \`a tour ses r\'ealisations,
    ses blocages, ses travaux en cours et ses prochaines \'etapes.
  \item[Revue des blocages]~: \`a chaque r\'eunion, les \'el\'ements
    bloquants sont identifi\'es et des solutions sont propos\'ees
    collectivement.
  \item[Validation des livraisons]~: \`a l'issue de chaque phase
    significative, l'\'equipe valide les fonctionnalit\'es livr\'ees
    avant de passer \`a la phase suivante.
\end{description}

\subsection{Outils collaboratifs}

L'\'equipe utilise un ensemble d'outils collaboratifs couvrant
l'ensemble des besoins du projet.

\begin{itemize}
  \item \textbf{GitHub}~: plateforme centrale du projet. Elle h\'eberge
    le code source, les issues, les pull requests, les discussions de
    code review et le tableau de bord de projet (GitHub Projects).
  \item \textbf{Jira}~: outil compl\'ementaire de suivi de projet pour
    la planification des t\^aches et le reporting d'avancement.
  \item \textbf{Discord}~: outil de communication quotidienne pour les
    \'echanges rapides, le partage d'\'ecran et l'organisation
    informelle.
  \item \textbf{Git}~: gestion de versions d\'ecentralis\'ee avec un
    workflow trunk-based. Les branches de fonctionnalit\'es sont
    cr\'e\'ees \`a partir de \texttt{main}, puis fusionn\'ees apr\`es
    validation via pull request.
  \item \textbf{Code review syst\'ematique}~: toute modification est
    soumise \`a une revue par au moins un autre membre de l'\'equipe
    avant d'\^etre int\'egr\'ee. Cette pratique garantit la qualit\'e
    du code et la diffusion des connaissances.
\end{itemize}

\section{Gestion des risques}

La gestion des risques est un aspect essentiel de la conduite du
projet. L'\'equipe a identifi\'e les risques principaux, \'evalu\'e
leur probabilit\'e et leur impact, et d\'efini des mesures de
mitigation. Le tableau~\ref{tab:risques} r\'esume cette analyse.

\begin{table}[h]
  \centering
  \caption{Analyse des risques du projet}
  \label{tab:risques}
  \begin{tabularx}{\textwidth}{@{}l c c X@{}}
    \toprule
    \textbf{Risque}
      & \textbf{Prob.}
      & \textbf{Impact}
      & \textbf{Mitigation} \\
    \midrule
    Indisponibilit\'e API externe (VT, Claude)
      & Moyenne
      & \'Elev\'e
      & D\'egradation gracieuse, fallback, cache \\
    Probl\`emes de compatibilit\'e multi-arch
      & Faible
      & Moyen
      & Tests CI sur les deux architectures \\
    D\'erive de planning
      & Moyenne
      & \'Elev\'e
      & Sprints courts, priorisation MVP \\
    S\'ecurit\'e des donn\'ees clients
      & Faible
      & Critique
      & Isolation r\'eseau, secrets, .env \\
    Mont\'ee en comp\'etence (Docker, CI/CD)
      & Moyenne
      & Moyen
      & Documentation, POCs, pair programming \\
    \bottomrule
  \end{tabularx}
\end{table}

Chaque risque fait l'objet d'un suivi r\'egulier pendant les sprints.
Les mesures de mitigation sont ajust\'ees si la probabilit\'e ou
l'impact \'evolue au cours du projet.

\section{Plan de d\'eploiement en phases}

Le projet LockBits est d\'ecompos\'e en six phases, chacune produisant
des livrables concrets. Ce phasage permet de r\'eduire les risques
d'int\'egration en validant chaque \'etape avant de passer \`a la
suivante.

\subsection{Phase 0~: Initialisation (Semaines 1--2)}

La phase d'initialisation pose les fondations du projet. Elle comprend
la mise en place du d\'ep\^ot GitHub principal avec ses
sous-modules (EDR, site), la configuration initiale de la pipeline
CI/CD, la cr\'eation des environnements de d\'eveloppement Docker, et
la d\'efinition des conventions de code et de commit.

\subsection{Phase 1~: Core EDR (Semaines 3--6)}

Cette phase est consacr\'ee au d\'eveloppement du c\^ocur du produit
EDR. Le serveur FastAPI est d\'evelopp\'e avec son API REST et sa base
SQLite. L'agent Python est cr\'e\'e avec les fonctionnalit\'es de scan,
de monitoring et de FIM. L'int\'egration VirusTotal est mise en place,
et les premiers tests unitaires sont \'écrits.

\subsection{Phase 2~: Site client (Semaines 7--10)}

Le portail client est d\'evelopp\'e durant cette phase. Le site
PHP/Apache est construit avec sa base de donn\'ees MySQL, le syst\`eme
d'authentification, le tableau de bord client et l'int\'egration GLPI
via SSO.

\subsection{Phase 3~: Chatbot (Semaines 11--12)}

Le chatbot est int\'egr\'e au site client. Cette phase inclut la
connexion \`a l'API Claude d'Anthropic, le d\'eveloppement de
l'interface de chat, la mise en place du rate limiting et les mesures
de s\'ecurit\'e associ\'ees.

\subsection{Phase 4~: CI/CD et d\'eploiement (Semaines 13--15)}

La pipeline CI/CD est finalis\'ee avec l'ensemble des \'etapes~: lint,
test, build multi-architecture, scan de s\'ecurit\'e et push vers
GHCR. Le script d'installation one-liner (\texttt{install.sh}) est
d\'evelopp\'e et test\'e. La documentation de d\'eploiement est
r\'edig\'ee.

\subsection{Phase 5~: Monitoring et s\'ecurit\'e (Semaines 16--18)}

La derni\`ere phase renforce la robustesse du syst\`eme. Prometheus,
Grafana et Loki sont d\'eploy\'es pour le monitoring. Les scans de
vuln\'erabilit\'es Trivy sont int\'egr\'es au pipeline. Le hardening
de la configuration est r\'ealis\'e et la documentation de s\'ecurit\'e
est finalis\'ee.

\section{Crit\`eres de succ\`es (MVP)}

Le Minimum Viable Product (MVP) est valid\'e lorsque les six crit\`eres
suivants sont atteints~:

\begin{enumerate}
  \item \textbf{Serveur EDR op\'erationnel} avec API REST document\'ee
    accessible via \texttt{/docs} (Swagger UI).
  \item \textbf{Agent EDR fonctionnel} capable de scan et de monitoring
    en temps r\'eel sur les endpoints.
  \item \textbf{Site client d\'eploy\'e} avec authentification, tableau
    de bord et gestion des clients.
  \item \textbf{Chatbot op\'erationnel} int\'egr\'e au site avec rate
    limiting fonctionnel.
  \item \textbf{Pipeline CI/CD verte} ex\'ecutant l'encha\^inement
    complète~: lint, test, build, push.
  \item \textbf{Installation en une commande} via le script
    \texttt{install.sh} utilisant \texttt{curl}.
\end{enumerate}

Ces crit\`eres garantissent que la solution est utilisable en
condition r\'eelle et peut \^etre d\'eploy\'ee par un client sans
intervention de l'\'equipe de d\'eveloppement.
